\documentclass[pgs]{DOE}

\begin{document}

\section{Background/Introduction (approximately 1 page)}

Explain the importance and relevance of the proposed work and clearly articulate which aspect of the chosen challenge will be addressed and solved. Set the proposed research in perspective to other efforts in the field. Highlight novel or unique aspects of the proposed work. Cite relevant literature.

\vspace{18pt}
\section{Project Objectives (approximately 0.5 page)}

Provide a clear and concise statement of the specific objectives of the proposed project. Address how the objectives align with the chosen focus topic and lead to an AI advantage.  

\vspace{18pt}
\section{Proposed Research and Methods (approximately 1.5 pages)}

Provide a clear research plan for a 9-month project and describe the proposed activities and methods. Include enough technical details to evaluate the impact of the proposed activity. For each activity, indicate the responsibility of the key investigator(s) and the associated budget. If the proposed application is building on an existing, currently funded DOE project, describe how the submitted application leverages that work and is distinct from the existing funding.

\vspace{18pt}
\section{Milestones in the Nine Months (approximately 1 page)}

Provide a list of clearly defined and measurable milestones of the project.  

\vspace{18pt}
\section{Data Sources and Models (approximately 0.5 page)}

Describe which existing or new data sources will be used. Address your team’s access to the data. If applicable, indicate which AI models or frameworks will be used in your proposed project. Do not include a description of computing resources and the means by which you will secure an allocation timed to your performance period; this should be done in Appendix 2.  

\vspace{18pt}
\section{Decision Gate Metrics (approximately 0.5 page)}

Provide specific metrics that can be used for the evaluation of Phase I go/no go decision. DOE encourages the development of metrics to identify AI advantage. One such metric could include scaling behavior which shows increasing performance as additional data, computing, and/or other resources are applied. Quantitative statements enabling transformative science and engineering through AI are preferred.

\end{document}